\chapter*{머리말}
\addcontentsline{toc}{chapter}{머리말}

Volume 1에서는 군, 환, 아이디얼과 다항식환을 배웠다. 이제 그 도구를 이용해 방정식의 근을 실제로 담는 새로운 체를 만들고, 서로 다른 체 사이의 크기와 구조를 비교한다. 이 과정이 \emph{체확장론}이며, 갈루아 이론의 직접적인 출발점이다.

이 권의 중심 질문은 다음과 같다.

\begin{quote}
다항식의 근을 첨가하면 어떤 체가 생기며, 그 체는 원래 체보다 얼마나 큰가?
\end{quote}

이를 위해 표수와 소체, 체확장의 차수, 최소다항식, 단순확장, 대수적 확장, 분해체, 분리성과 정상성을 차례로 다룬다. 최종적으로는 ``체의 중간구조''와 ``자동동형군의 부분군''이 서로 대응한다는 갈루아 이론의 기본정리에 도달할 준비를 갖춘다.

\section*{권장 선수지식}

다음 내용에 익숙하다고 가정한다.

\begin{itemize}
  \item 벡터공간의 기저와 차원
  \item 환 준동형, 아이디얼과 몫환
  \item 체 위 다항식환의 나눗셈 알고리즘
  \item 기약다항식과 $K[X]/(f)$의 체 구조
  \item 군 준동형과 몫군의 기본 개념
\end{itemize}

필요한 결과는 사용 직전에 짧게 복습하므로 Volume 1의 모든 증명을 암기할 필요는 없다.

\section*{학습 방법}

각 절은 동기, 정의, 예제, 정리, 증명 전략, 완전한 증명, 이해 점검의 순서로 구성한다. 처음 읽을 때에는 정의와 핵심 예제를 우선 이해하고, 두 번째 읽기에서 증명을 직접 재구성하는 편이 좋다.

연습문제는 A, B, C의 세 단계로 나눈다.

\begin{itemize}
  \item A: 정의 확인과 계산
  \item B: 핵심 정리의 재증명
  \item C: 여러 개념을 연결하는 구조 문제
\end{itemize}

별표 문제는 뒤의 선별 풀이에서 상세히 다룬다. 풀이를 보기 전 최소한 정의와 사용할 정리를 직접 적어 보는 것을 권한다.

\section*{현재 판의 범위}

Draft 0.1에는 체확장의 첫 장인 표수, 소체, 차수, 탑 법칙, 대수적 원소, 최소다항식과 유한생성 대수적 확장을 수록한다. 이후 장에서 대수적 폐포, 분해체, 분리성, 순수 비분리 확장과 유한체를 추가한다.
